% !TEX root = master.tex
% !TEX program = lualatex

\chapter{Pointers - 2}\label{chap:pointers_} % (fold)

\section{C01b — Pointers : Introduction à l’allocation dynamique}

\subsection{Rappel : les trois usages des pointeurs}

Lors du premier cours sur les pointeurs, nous avons vu trois usages principaux :

\begin{enumerate}
    \item \textbf{Passage par référence}  
    Permettre à une fonction de modifier une variable définie à l’extérieur.

    \item \textbf{Description générique d’objets}  
    Manipuler des objets dont la nature exacte n’est pas connue au moment de l’écriture du code.

    \item \textbf{Allocation dynamique}  
    Gérer nous-mêmes la durée de vie d’une zone mémoire.
\end{enumerate}

Les deux premiers usages ont été détaillés précédemment.  
Nous abordons maintenant le troisième : \textbf{l’allocation dynamique}.

\subsection{Gestion de la mémoire d’un processus}

Pour comprendre l’allocation dynamique, il faut d’abord comprendre 
\textbf{l’organisation mémoire d’un programme}.

La mémoire virtuelle d’un processus peut être schématisée ainsi :

\begin{itemize}
    \item En bas : le \textbf{code du programme} + constantes.
    \item Puis : les \textbf{variables globales et statiques}.
    \item Ensuite : le \textbf{tas (heap)}.
    \item En haut : la \textbf{pile (stack)}.
\end{itemize}

La pile et le tas évoluent en sens opposés :

\begin{itemize}
    \item La \textbf{pile} grandit vers les adresses décroissantes.
    \item Le \textbf{tas} grandit vers les adresses croissantes.
\end{itemize}

\subsection{Allocation statique (automatique)}

Jusqu’à présent, toutes les variables déclarées dans vos programmes C
étaient allouées automatiquement sur la \textbf{pile}.

Exemple :

\begin{lstlisting}[language=C]
void f() {
    int i = 5;
}
\end{lstlisting}

Lorsque la fonction commence :
\begin{itemize}
    \item de la mémoire est réservée pour \texttt{i} sur la pile.
\end{itemize}

Lorsque la fonction se termine :
\begin{itemize}
    \item la mémoire est automatiquement libérée.
\end{itemize}

On parle d’\textbf{allocation automatique (statique du point de vue du programmeur)}.

Le programmeur ne contrôle pas la durée de vie exacte :
elle correspond strictement à la portée (\textit{scope}) de la variable.

\subsection{Limite de la pile}

La pile a une taille limitée (souvent quelques mégaoctets).

Sous Unix/Linux, on peut vérifier la taille maximale avec :

\begin{lstlisting}[language=C]
ulimit -s
\end{lstlisting}

Si l’on alloue des structures trop volumineuses sur la pile,
on risque un \textbf{stack overflow}.

\subsection{Allocation dynamique (sur le tas)}

Lorsque l’on veut :

\begin{enumerate}
    \item Contrôler la \textbf{durée de vie} d’un objet au-delà de la portée,
    \item Allouer de \textbf{grandes quantités de mémoire},
\end{enumerate}

il faut utiliser le \textbf{tas (heap)}.

Contrairement à la pile :

\begin{itemize}
    \item L’allocation est faite explicitement par le programmeur.
    \item La désallocation doit aussi être faite explicitement.
    \item La mémoire disponible est beaucoup plus grande.
\end{itemize}

Le tas est limité uniquement par la mémoire virtuelle disponible 
(et bien sûr la mémoire physique réelle).

\subsection{Idée fondamentale : durée de vie}

La notion clé ici est la \textbf{durée de vie d’un objet en mémoire}.

\begin{itemize}
    \item Sur la pile : durée de vie = durée de la portée.
    \item Sur le tas : durée de vie = décidée par le programmeur.
\end{itemize}

L’allocation dynamique permet donc de dissocier :

\[
\text{durée de vie} \neq \text{durée de la portée}
\]

C’est précisément le troisième grand usage des pointeurs.


\section{C01b — Pointers : malloc et calloc}
\subsection{Allocation dynamique : \texttt{malloc}, \texttt{calloc}, \texttt{realloc} et \texttt{free}}

\subsubsection{Cycle de vie d’une allocation dynamique}

Lorsque l’on décide de gérer dynamiquement une zone mémoire, son cycle de vie est le suivant :

\begin{center}
Allocation $\longrightarrow$ Utilisation $\longrightarrow$ Libération
\end{center}

Les fonctions principales sont :

\begin{itemize}
    \item \texttt{malloc}
    \item \texttt{calloc}
    \item \texttt{realloc}
    \item \texttt{free}
\end{itemize}

\subsubsection{\texttt{malloc}}

Prototype :

\begin{lstlisting}[language=C]
void* malloc(size_t size);
\end{lstlisting}

\texttt{malloc} réserve \texttt{size} octets sur le tas et retourne un pointeur
vers le début de la zone allouée.

Exemple :

\begin{lstlisting}[language=C]
double* p = malloc(sizeof(double));
\end{lstlisting}

On utilise généralement l’opérateur :

\begin{lstlisting}[language=C]
sizeof(type)
\end{lstlisting}

qui retourne la taille en octets du type donné.

Important :

\begin{itemize}
    \item La mémoire retournée par \texttt{malloc} n’est PAS initialisée.
    \item Son contenu est indéterminé.
\end{itemize}

\subsubsection{\texttt{calloc}}

Prototype :

\begin{lstlisting}[language=C]
void* calloc(size_t nmemb, size_t size);
\end{lstlisting}

\texttt{calloc} alloue :

\[
nmemb \times size
\]

octets consécutifs en mémoire.

Exemple :

\begin{lstlisting}[language=C]
double* p = calloc(3, sizeof(double));
\end{lstlisting}

Cela alloue la place pour 3 doubles consécutifs.

Différences avec \texttt{malloc} :

\begin{itemize}
    \item \texttt{calloc} protège contre les erreurs d’overflow
          dans la multiplication.
    \item \texttt{calloc} initialise la mémoire à 0.
\end{itemize}

\subsubsection{Pourquoi préférer \texttt{calloc} ? (Bug OpenSSH)}

Considérons le code suivant (simplifié) :

\begin{lstlisting}[language=C]
int nresp = get_number();
char** response;

if (nresp > 0) {
    response = malloc(nresp * sizeof(char*));
    for (int i = 0; i < nresp; i++) {
        response[i] = get_packet();
    }
}
\end{lstlisting}

Supposons :

\[
nresp = 1\,073\,741\,856
\]

Sur une architecture 32 bits :

\begin{itemize}
    \item \texttt{sizeof(char*) = 4}
    \item Donc on calcule :
    \[
    1\,073\,741\,856 \times 4
    \]
\end{itemize}

Cette valeur dépasse la capacité représentable sur 32 bits.
Elle déborde et devient 128.

Résultat :

\begin{itemize}
    \item On alloue 128 octets seulement.
    \item La boucle itère jusqu’à 1 milliard.
    \item On écrit largement hors de la zone allouée.
\end{itemize}

C’est un \textbf{buffer overflow}, faille de sécurité réelle observée dans OpenSSH.

Solution :

\begin{lstlisting}[language=C]
response = calloc(nresp, sizeof(char*));
\end{lstlisting}

\texttt{calloc} détecte l’overflow interne et évite ce problème.

\subsubsection{Initialisation mémoire : \texttt{memset}}

Si l’on utilise \texttt{malloc} mais que l’on souhaite initialiser :

\begin{lstlisting}[language=C]
#include <string.h>

memset(ptr, value, size);
\end{lstlisting}

Exemple :

\begin{lstlisting}[language=C]
int* p = malloc(sizeof(int));
memset(p, 0, sizeof(int));
\end{lstlisting}

Cela met tous les octets à 0.

Important :

\begin{itemize}
    \item \texttt{memset} agit au niveau des octets.
    \item Mettre tous les bits à 0 correspond à la valeur 0
          pour les types entiers et les pointeurs nuls.
\end{itemize}

\subsubsection{Vérification obligatoire}

\texttt{malloc} et \texttt{calloc} retournent \texttt{NULL}
si l’allocation échoue.

Toujours écrire :

\begin{lstlisting}[language=C]
double* p = calloc(3, sizeof(double));
if (p == NULL) {
    // gérer l'erreur
    exit(EXIT_FAILURE);
}
\end{lstlisting}

Ne jamais utiliser un pointeur sans avoir vérifié qu’il n’est pas nul.

\subsubsection{\texttt{free}}

Lorsque la zone mémoire n’est plus nécessaire :

\begin{lstlisting}[language=C]
free(p);
\end{lstlisting}

Après un \texttt{free}, le pointeur devient invalide.
Bonne pratique :

\begin{lstlisting}[language=C]
free(p);
p = NULL;
\end{lstlisting}

\subsection*{Résumé}

\begin{center}
\begin{tabular}{|c|c|c|}
\hline
Fonction & Initialise ? & Protège overflow ? \\
\hline
malloc & Non & Non \\
calloc & Oui (0) & Oui \\
\hline
\end{tabular}
\end{center}

Principes essentiels :

\begin{itemize}
    \item Toujours utiliser \texttt{sizeof}.
    \item Toujours vérifier le retour.
    \item Toujours libérer avec \texttt{free}.
    \item Préférer \texttt{calloc} si possible.
\end{itemize}


\section{C01b — Pointers : realloc}

\subsection{Libération de la mémoire : \texttt{free}}

\subsubsection{Fin de la durée de vie}

Dans le cycle de vie d’une allocation dynamique :

\[
\text{Allocation} \longrightarrow \text{Utilisation} \longrightarrow \text{Libération}
\]

après avoir utilisé une zone mémoire allouée avec
\texttt{malloc}, \texttt{calloc} ou \texttt{realloc},
il faut la rendre au système.

C’est le rôle de :

\begin{lstlisting}[language=C]
free(ptr);
\end{lstlisting}

Ne pas libérer la mémoire entraîne une \textbf{fuite mémoire}
(\textit{memory leak}) :

\begin{itemize}
    \item Le programme consomme de plus en plus de mémoire.
    \item La mémoire n’est jamais rendue au système.
    \item Sur le long terme : ralentissement, crash.
\end{itemize}

\subsubsection{Ce que fait réellement \texttt{free}}

Attention à un abus de langage courant :

\begin{center}
\textit{On ne libère pas le pointeur.}
\end{center}

On libère \textbf{la zone mémoire pointée}.

Schéma conceptuel :

\begin{itemize}
    \item Monde des adresses : le pointeur contient une adresse.
    \item Monde des valeurs : la zone mémoire contient des données.
\end{itemize}

Lorsque l’on fait :

\begin{lstlisting}[language=C]
free(ptr);
\end{lstlisting}

\begin{itemize}
    \item La zone mémoire est rendue au système.
    \item Mais le pointeur contient toujours la même adresse.
\end{itemize}

Le pointeur devient alors un \textbf{dangling pointer}
(pointeur pendu).

Utiliser ce pointeur après un \texttt{free}
entraîne un comportement indéfini.

\subsubsection{Bonne pratique essentielle}

Toujours écrire :

\begin{lstlisting}[language=C]
free(ptr);
ptr = NULL;
\end{lstlisting}

Pourquoi ?

\begin{itemize}
    \item On rend la mémoire.
    \item On efface l’adresse.
    \item On évite toute utilisation accidentelle.
\end{itemize}

Un pointeur nul est sûr : toute tentative d’accès
provoquera une erreur détectable.

\subsubsection{Règle fondamentale}

\begin{center}
\textbf{Toute allocation dynamique doit avoir un \texttt{free} correspondant.}
\end{center}

Exemples :

\begin{itemize}
    \item 1 appel à \texttt{malloc} $\rightarrow$ 1 appel à \texttt{free}
    \item 25 appels à \texttt{calloc} dans une boucle $\rightarrow$
          25 appels à \texttt{free}
\end{itemize}

Sinon : fuite mémoire.

\subsubsection{Erreur classique}

\begin{lstlisting}[language=C]
int* p = malloc(sizeof(int));
...
free(p);
*p = 5;   // ERREUR : accès après free
\end{lstlisting}

C’est un accès à une zone mémoire que l’on n’a plus le droit d’utiliser.

\subsection*{Résumé}

\begin{itemize}
    \item \texttt{free} libère la zone pointée, pas le pointeur.
    \item Toujours mettre le pointeur à \texttt{NULL} après.
    \item Toute allocation dynamique doit être libérée.
    \item Les fuites mémoire sont des erreurs graves en programmation système.
\end{itemize}

\section{C01b — Pointers : Allocation dynamique : exemple complet avec \texttt{malloc} et \texttt{free}}

Nous considérons ici le cas où l’on souhaite créer un entier dont la durée de vie dépasse
la portée d’une fonction. Il s’agit du cas classique où une fonction alloue dynamiquement
une zone mémoire et retourne un pointeur vers celle-ci.

L’objectif est de comprendre clairement la différence entre :
\begin{itemize}
    \item la durée de vie des \textbf{pointeurs} (variables locales),
    \item la durée de vie de la \textbf{zone mémoire pointée}.
\end{itemize}

\subsection{1. Fonction de création}

La fonction suivante crée dynamiquement un entier et retourne son adresse.

\begin{lstlisting}[language=C]
#include <stdio.h>
#include <stdlib.h>

int* createLongLiveInt(void) {
    int* px = NULL;              // px existe uniquement dans cette fonction

    px = malloc(sizeof(int));    // allocation dynamique

    if (px != NULL) {
        *px = 20;                // initialisation de la zone mémoire
    }

    return px;                   // on retourne l'adresse allouee
}
\end{lstlisting}

\paragraph{Analyse détaillée}

\begin{itemize}
    \item \texttt{px} est un pointeur local : il n'existe que dans \texttt{createLongLiveInt}.
    \item \texttt{malloc(sizeof(int))} alloue une zone mémoire capable de contenir un entier.
    \item \texttt{malloc} retourne l'adresse de cette zone.
    \item Cette adresse est copiée dans \texttt{px}.
    \item On écrit ensuite dans la zone pointée via \texttt{*px = 20}.
    \item Lors du \texttt{return px}, on retourne une copie de l'adresse.
\end{itemize}

Important : à la fin de la fonction, la variable \texttt{px} disparaît,
mais \textbf{la zone mémoire allouée continue d'exister} tant qu'elle n'est pas libérée avec \texttt{free}.

\subsection{2. Utilisation dans une autre fonction}

Voyons maintenant comment cette zone peut être utilisée ailleurs :

\begin{lstlisting}[language=C]
void exampleUsage(void) {
    int* q = createLongLiveInt();   // q recoit l'adresse

    if (q != NULL) {
        *q = 123;                   // modification de la meme zone
    }

    int* r = q;                     // copie d'adresse (egalite de pointeurs)

    printf("%d\n", *r);             // acces via r

    free(q);                        // liberation de la zone
    q = NULL;                       // bonne pratique
    r = NULL;                       // eviter pointeur pendu
}
\end{lstlisting}

\subsection{3. Égalité de pointeurs}

L’instruction :

\[
r = q;
\]

ne copie pas la valeur stockée dans la mémoire pointée,
mais uniquement l'adresse.

Ainsi :
\begin{itemize}
    \item \texttt{q} et \texttt{r} pointent vers la même zone mémoire.
    \item Modifier \texttt{*q} modifie ce que voit \texttt{r}.
\end{itemize}

Il s'agit d'une \textbf{copie d'adresse}, pas d'une copie de contenu.

\subsection{4. Durée de vie : ce qu’il faut absolument comprendre}

La gestion de la durée de vie concerne :

\[
\textbf{la zone mémoire allouée}, \quad \textbf{pas les pointeurs}.
\]

\begin{itemize}
    \item \texttt{px} disparaît à la fin de \texttt{createLongLiveInt}.
    \item \texttt{q} disparaît à la fin de \texttt{exampleUsage}.
    \item Mais la zone mémoire reste valide tant qu'on n'a pas fait \texttt{free}.
\end{itemize}

La séquence correcte est donc :

\[
\texttt{malloc} \rightarrow \text{utilisation} \rightarrow \texttt{free}
\]

C’est le programmeur qui décide explicitement de la durée de vie.

\subsection{5. Attention aux erreurs classiques}

\paragraph{1. Oublier le \texttt{free}}

Provoque une fuite mémoire (memory leak).

\paragraph{2. Pointeur = NULL après \texttt{free}}
\\
Après :
\begin{lstlisting}[language=C]
free(q);
\end{lstlisting}

la zone mémoire est libérée, mais l’adresse contenue dans \texttt{q}
reste inchangée. Le pointeur devient un \textbf{dangling pointer}.

Toujours faire :

\begin{lstlisting}[language=C]
q = NULL;
\end{lstlisting}

\subsection{6. Schéma conceptuel temporel}

On peut résumer ainsi :

\begin{itemize}
    \item Allocation dans \texttt{createLongLiveInt}
    \item Transmission de l’adresse par copie
    \item Utilisation via plusieurs pointeurs
    \item Libération explicite avec \texttt{free}
\end{itemize}

La zone mémoire existe indépendamment des variables locales
qui ont servi à transmettre son adresse.

\bigskip
\textbf{Conclusion :}

L’allocation dynamique permet de créer des objets dont la durée de vie
n’est pas limitée à la portée d’une fonction.
Elle donne un contrôle total au programmeur,
mais impose une discipline stricte :

\[
\boxed{\text{Chaque malloc doit correspondre à un free}}
\]

\section{C01b — Pointers : Initialisation des pointeurs et bonnes pratiques}

Un des pièges classiques avec les pointeurs est l’utilisation de zones mémoire
non valides, en particulier des pointeurs non initialisés.

\subsection{1. Le cas dangereux : pointeur non initialisé}

Considérons le code suivant :

\begin{lstlisting}[language=C]
#include <stdio.h>

void badExample(void) {
    int* px;          // NON initialise

    // ... beaucoup de code ici ...

    *px = 20;         // Ecriture dans une zone inconnue
}
\end{lstlisting}

\paragraph{Que se passe-t-il ici ?}

\begin{itemize}
    \item \texttt{px} est déclaré mais jamais initialisé.
    \item Il contient donc une valeur indéterminée.
    \item L'instruction \texttt{*px = 20;} tente d'écrire à une adresse inconnue.
\end{itemize}

Le programme compile parfaitement (aucune erreur syntaxique),
mais à l’exécution il y a de fortes chances d’obtenir :

\[
\textbf{Segmentation fault}
\]

Pourquoi ?

Parce que vous accédez à une zone mémoire à laquelle vous n’avez pas le droit
d’accéder.

\subsection{2. Segmentation fault vs Bus error}

\paragraph{Segmentation fault}

Se produit lorsque :
\begin{itemize}
    \item vous accédez à une zone mémoire interdite,
    \item le noyau détecte directement l'accès invalide.
\end{itemize}

\paragraph{Bus error}

Plus rare. Peut apparaître lorsque :
\begin{itemize}
    \item l’adresse pointe vers une zone mal alignée,
    \item ou vers de la mémoire matérielle spécifique (périphérique),
    \item et que l’erreur est détectée au niveau du bus matériel.
\end{itemize}

Dans les deux cas : c’est un accès mémoire invalide.

\subsection{3. Bonne pratique : toujours initialiser}

Il est fortement recommandé d’écrire :

\begin{lstlisting}[language=C]
int* px = NULL;
\end{lstlisting}

Pourquoi ?

Cela ne supprime pas le bug si vous oubliez d’allouer,
mais garantit que le programme plantera exactement à l’endroit
où vous utilisez le pointeur.

Exemple :

\begin{lstlisting}[language=C]
int* px = NULL;
*px = 20;   // Segmentation fault immediate
\end{lstlisting}

Cela rend le débogage beaucoup plus simple.

\subsection{4. Règles fondamentales de gestion mémoire}

\paragraph{Règle 1 : Toute allocation doit être libérée}

\[
\boxed{\texttt{malloc} \longleftrightarrow \texttt{free}}
\]

Exemple correct :

\begin{lstlisting}[language=C]
int* px = malloc(sizeof(int));
if (px != NULL) {
    *px = 10;
}
free(px);
px = NULL;
\end{lstlisting}

Sinon : fuite mémoire (\textit{memory leak}).

\subsection{5. Responsabilité : qui alloue doit libérer}

Principe fondamental :

\begin{itemize}
    \item Si une fonction alloue, elle doit documenter clairement
    qui est responsable du \texttt{free}.
    \item Idéalement : fournir une fonction de destruction.
\end{itemize}

Exemple conceptuel :

\begin{lstlisting}[language=C]
Matrix* createMatrix(int n);
void destroyMatrix(Matrix* m);
\end{lstlisting}

Toujours proposer un couple constructeur / destructeur.

\subsection{6. Toujours tester le retour de malloc / calloc}

\begin{lstlisting}[language=C]
int* px = malloc(sizeof(int));

if (px == NULL) {
    fprintf(stderr, "Allocation failed\n");
    exit(EXIT_FAILURE);
}
\end{lstlisting}

Ne jamais supposer que l’allocation réussit.

\subsection{7. malloc vs calloc}

Pour des allocations multiples :

\begin{lstlisting}[language=C]
int* array = calloc(n, sizeof(int));
\end{lstlisting}

Avantages :

\begin{itemize}
    \item Initialise la mémoire à 0
    \item Evite certaines erreurs d’initialisation
\end{itemize}

Toujours écrire :

\begin{lstlisting}[language=C]
int* array = malloc(n * sizeof(int));
\end{lstlisting}

et jamais :

\begin{lstlisting}[language=C]
int* array = malloc(n * 4);   // Mauvaise pratique
\end{lstlisting}

\subsection{8. Initialiser la mémoire}

Deux solutions :

\begin{itemize}
    \item Utiliser \texttt{calloc}
    \item Utiliser \texttt{memset}: void *memset(void *adresseDeDepart, int valeurACopier, nbOctetARemplir n);
\end{itemize}

\begin{lstlisting}[language=C]
int* array = malloc(n * sizeof(int));
memset(array, 0, n * sizeof(int));
\end{lstlisting}

\subsection{9. Après free : effacer l’adresse}

\begin{lstlisting}[language=C]
free(px);
px = NULL;
\end{lstlisting}

Sinon, \texttt{px} devient un \textbf{dangling pointer}
(pointeur pendu).

\subsection{10. Passage par pointeur et const}

Si un pointeur est utilisé uniquement en lecture :

\begin{lstlisting}[language=C]
void printValue(const int* px);
\end{lstlisting}

Cela garantit que la fonction ne modifiera pas la zone pointée.

\textbf{Résumé des bonnes pratiques}

\begin{itemize}
    \item Toujours initialiser les pointeurs.
    \item Toujours tester les allocations.
    \item Toujours libérer ce qui a été alloué.
    \item Mettre le pointeur à \texttt{NULL} après \texttt{free}.
    \item Utiliser \texttt{sizeof(type)}.
    \item Fournir des fonctions de destruction.
    \item Utiliser \texttt{const} quand c’est possible.
\end{itemize}

\[
\boxed{\text{La mémoire dynamique donne du pouvoir, mais exige de la discipline.}}
\]

\section{C01b — Pointers : Étude de cas : implémentation d’un tableau dynamique (vector)}

Nous commençons maintenant une étude pratique visant à implémenter
un tableau dynamique en C, c’est-à-dire l’équivalent :

\begin{itemize}
    \item d’un \texttt{ArrayList} en Java,
    \item d’un \texttt{std::vector} en C++.
\end{itemize}

Le langage C ne fournit pas ce type de structure dans sa bibliothèque standard.
Nous allons donc l’implémenter nous-mêmes, étape par étape.

\subsection*{1. Objectif}

Nous voulons créer un tableau :

\begin{itemize}
    \item dont la taille n’est pas connue à l’avance,
    \item et qui peut grandir dynamiquement.
\end{itemize}

Cela implique :
\begin{itemize}
    \item une allocation dynamique,
    \item une réallocation lorsque la capacité est dépassée,
    \item une libération explicite.
\end{itemize}

\subsection*{2. Interface conceptuelle}

On souhaite pouvoir écrire quelque chose comme :

\begin{lstlisting}[language=C]
vector v;

if (vector_create(&v) == 0) {
    vector_push(&v, 10);
    vector_push(&v, 42);
    vector_destroy(&v);
}
\end{lstlisting}

Remarques importantes :

\begin{itemize}
    \item Les fonctions reçoivent un pointeur sur \texttt{vector}
    car elles modifient la structure.
    \item Si on alloue, on doit fournir une fonction de destruction.
\end{itemize}

\subsection*{3. Définition du type des éléments}

Pour rendre l’implémentation flexible, on définit :

\begin{lstlisting}[language=C]
typedef int element_type;
\end{lstlisting}

En changeant une seule ligne, on pourrait transformer notre
vector d’entiers en vector de \texttt{double}, par exemple.

\subsection*{4. Structure interne du vector}

Nous définissons maintenant la structure :

\begin{lstlisting}[language=C]
typedef struct {
    size_t size;        // nombre d'elements effectivement stockes
    size_t capacity;    // nombre maximal d'elements alloues
    element_type* data; // pointeur vers la zone memoire
} vector;
\end{lstlisting}

Il est crucial de comprendre la différence entre :

\begin{itemize}
    \item \texttt{size} : nombre réel d’éléments utilisés,
    \item \texttt{capacity} : nombre maximal d’éléments pouvant être stockés
    sans nouvelle allocation.
\end{itemize}

\subsection*{5. Sémantique mémoire}

Supposons :

\[
\texttt{capacity} = 100
\]

Si un \texttt{int} occupe 4 octets, alors la zone allouée correspond à :

\[
100 \times \texttt{sizeof(int)} = 400 \text{ octets}
\]

Visuellement :

\[
\underbrace{[\,\square\,][\,\square\,]\dots[\,\square\,]}_{\text{capacity cases}}
\]

Mais si seulement 5 éléments sont utilisés :

\[
\texttt{size} = 5
\]

Les autres cases sont simplement disponibles pour de futurs ajouts.

\subsection*{6. Stratégie d’extension}

Pourquoi stocker une capacité plus grande que la taille réelle ?

Parce que si on réalloue à chaque insertion :

\begin{lstlisting}[language=C]
push -> realloc -> push -> realloc -> push -> realloc
\end{lstlisting}

cela serait très coûteux.

On adopte donc une stratégie par blocs :

\begin{itemize}
    \item Allouer par exemple 100 éléments.
    \item Quand la capacité est atteinte, doubler la capacité.
\end{itemize}

\[
100 \rightarrow 200 \rightarrow 400 \rightarrow 800
\]

Cela permet d’avoir une complexité amortie efficace.

\subsection*{7. Problème : que faire au 101\textsuperscript{e} élément ?}

Supposons :

\begin{itemize}
    \item \texttt{capacity = 100}
    \item \texttt{size = 100}
\end{itemize}

Lorsqu’on veut insérer le 101\textsuperscript{e} élément,
il faut agrandir la zone mémoire.

C’est ici qu’intervient :

\[
\boxed{\texttt{realloc}}
\]

\subsection*{8. Rôle de realloc}

\texttt{realloc} permet de :

\begin{itemize}
    \item agrandir ou réduire une zone mémoire existante,
    \item éventuellement déplacer la zone en mémoire,
    \item retourner une nouvelle adresse si nécessaire.
\end{itemize}

Principe général :

\begin{lstlisting}[language=C]
pointer = realloc(pointer, nouvelle_taille);
\end{lstlisting}

Attention :

\begin{itemize}
    \item Si \texttt{realloc} échoue, il retourne \texttt{NULL}.
    \item L’ancien pointeur reste valide.
    \item Il faut donc toujours tester le retour.
\end{itemize}

\subsection*{9. Résumé conceptuel}

Un vector dynamique contient :

\begin{itemize}
    \item un nombre d’éléments utilisés,
    \item une capacité allouée,
    \item un pointeur vers la mémoire.
\end{itemize}

Lorsqu’on dépasse la capacité :

\begin{itemize}
    \item on appelle \texttt{realloc},
    \item on met à jour \texttt{capacity},
    \item on continue l’insertion.
\end{itemize}

\bigskip
\textbf{Idée clé :}

La structure elle-même est statique,
mais la mémoire pointée est dynamique et peut être déplacée.

\[
\boxed{\text{Le vector contrôle dynamiquement sa propre zone mémoire.}}
\]

\section{C01b — Pointers : La fonction \texttt{realloc}}

Pour redimensionner une zone mémoire déjà allouée dynamiquement,
le langage C fournit la fonction :

\[
\boxed{\texttt{realloc}}
\]

Attention : ce n’est pas un \texttt{recalloc}.
C’est un \textbf{re-malloc}.  
La nouvelle zone ajoutée n’est pas initialisée.

\subsection*{1. Prototype}

\begin{lstlisting}[language=C]
void* realloc(void* ptr, size_t new_size);
\end{lstlisting}

Arguments :

\begin{itemize}
    \item \texttt{ptr} : ancien pointeur (zone déjà allouée)
    \item \texttt{new\_size} : nouvelle taille en octets
\end{itemize}

Valeur de retour :

\begin{itemize}
    \item Nouvelle adresse (identique ou différente)
    \item \texttt{NULL} si échec
\end{itemize}

\subsection*{2. Comportement général}

\paragraph{Cas 1 : agrandissement}

Si on augmente la taille :

\begin{lstlisting}[language=C]
int* tmp = realloc(ptr, 200 * sizeof(int));
\end{lstlisting}

Alors :

\begin{itemize}
    \item Si de l’espace est disponible après la zone actuelle,
          l’adresse peut rester identique.
    \item Sinon, une nouvelle zone est allouée ailleurs.
    \item Les anciennes données sont copiées.
    \item L’ancienne zone est libérée automatiquement.
\end{itemize}

Important :
la partie nouvellement allouée n’est pas initialisée.

\subsection*{3. Cas 2 : réduction}

On peut aussi réduire la taille :

\begin{lstlisting}[language=C]
ptr = realloc(ptr, 50 * sizeof(int));
\end{lstlisting}

Dans ce cas :

\begin{itemize}
    \item La zone peut être conservée telle quelle.
    \item Ou éventuellement déplacée.
    \item Les données au-delà de la nouvelle taille sont perdues.
\end{itemize}

\subsection*{4. Cas 3 : échec}

Si \texttt{realloc} échoue :

\begin{itemize}
    \item Il retourne \texttt{NULL}.
    \item L’ancienne zone mémoire reste inchangée.
\end{itemize}

\paragraph{Mauvaise pratique}

\begin{lstlisting}[language=C]
ptr = realloc(ptr, new_size);   // DANGEREUX
\end{lstlisting}

Si \texttt{realloc} retourne \texttt{NULL},
vous perdez l’adresse originale → fuite mémoire.

\paragraph{Bonne pratique}

\begin{lstlisting}[language=C]
int* tmp = realloc(ptr, new_size);

if (tmp == NULL) {
    // echec : ptr est toujours valide
    fprintf(stderr, "Reallocation failed\n");
} else {
    ptr = tmp;
}
\end{lstlisting}

Toujours utiliser un pointeur temporaire.

\subsection*{5. Cas particuliers}

\paragraph{realloc(NULL, size)}

\begin{lstlisting}[language=C]
ptr = realloc(NULL, size);
\end{lstlisting}

Equivalent à :

\begin{lstlisting}[language=C]
ptr = malloc(size);
\end{lstlisting}

Même si c’est autorisé, il est préférable
d’utiliser explicitement \texttt{malloc}.

\paragraph{realloc(ptr, 0)}

\begin{lstlisting}[language=C]
realloc(ptr, 0);
\end{lstlisting}

Equivalent à :

\begin{lstlisting}[language=C]
free(ptr);
\end{lstlisting}

Là aussi, il est recommandé d’utiliser directement \texttt{free}.

\subsection*{6. Résumé du fonctionnement interne}

\begin{itemize}
    \item Peut agrandir ou réduire une zone mémoire.
    \item Peut déplacer la zone si nécessaire.
    \item Copie automatiquement les données existantes.
    \item Libère l’ancienne zone si déplacement.
    \item Ne touche pas à l’ancienne zone en cas d’échec.
    \item N’initialise pas la nouvelle mémoire ajoutée.
\end{itemize}

\bigskip
\textbf{Point essentiel :}

\[
\boxed{\texttt{realloc} = \text{redimensionnement intelligent + copie automatique}}
\]

Mais :

\[
\boxed{\text{La mémoire ajoutée n’est PAS initialisée}}
\]

Donc si vous agrandissez un tableau,
vous devez initialiser manuellement la nouvelle partie.

\section{C01b — Pointers : Étude de cas : implémentation complète du vector (constructeur, push, enlarge)}

Nous détaillons maintenant les fonctions principales :
\begin{itemize}
    \item le constructeur,
    \item le destructeur,
    \item la fonction \texttt{vector\_push},
    \item la fonction interne \texttt{vector\_enlarge}.
\end{itemize}

L’objectif est de comprendre :
\begin{itemize}
    \item le passage par référence du vector lui-même (cas 1),
    \item la gestion dynamique de la zone pointée (cas 3),
    \item l’écriture dite « atomique » pour éviter les corruptions en cas d’échec.
\end{itemize}

\subsection{1. Constantes et types}

\begin{lstlisting}[language=C]
#include <stdlib.h>
#include <stdio.h>
#include <limits.h>

typedef int element_type;

#define VECTOR_PADDING 128

typedef struct {
    size_t size;        // nombre d'elements utilises
    size_t allocated;   // capacite allouee
    element_type* content; // zone memoire dynamique
} vector;
\end{lstlisting}

\subsection{2. Constructeur}

Convention :
\begin{itemize}
    \item retourne \texttt{NULL} en cas d’échec,
    \item sinon retourne le pointeur reçu (pour chaînage possible).
\end{itemize}

\begin{lstlisting}[language=C]
vector* vector_create(vector* v) {
    if (v == NULL) {
        return NULL;
    }

    vector result;
    result.size = 0;
    result.allocated = VECTOR_PADDING;
    result.content = calloc(VECTOR_PADDING, sizeof(element_type));

    if (result.content == NULL) {
        return NULL;
    }

    *v = result;  // ecriture atomique finale
    return v;
}
\end{lstlisting}

Pourquoi travailler sur \texttt{result} ?

Pour ne modifier \texttt{v} que lorsque tout s’est bien passé.

\subsection{3. Destructeur}

\begin{lstlisting}[language=C]
void vector_destroy(vector* v) {
    if (v == NULL) {
        return;
    }

    free(v->content);
    v->content = NULL;
    v->size = 0;
    v->allocated = 0;
}
\end{lstlisting}

\subsection{4. Fonction interne : vector\_enlarge}

Objectif :
augmenter la capacité de \texttt{VECTOR\_PADDING} éléments.

Convention :
\begin{itemize}
    \item retourne \texttt{NULL} en cas d’échec,
    \item sinon retourne \texttt{v}.
\end{itemize}

Protection supplémentaire :
éviter l’overflow lors de la multiplication
\[
allocated \times sizeof(element\_type)
\]

\begin{lstlisting}[language=C]
vector* vector_enlarge(vector* v) {
    if (v == NULL) {
        return NULL;
    }

    vector result = *v;

    // nouvelle capacite
    result.allocated += VECTOR_PADDING;

    // protection contre overflow
    if (result.allocated > SIZE_MAX / sizeof(element_type)) {
        return NULL;
    }

    element_type* tmp = realloc(
        result.content,
        result.allocated * sizeof(element_type)
    );

    if (tmp == NULL) {
        return NULL;
    }

    result.content = tmp;

    *v = result;   // ecriture atomique
    return v;
}
\end{lstlisting}

Pourquoi tester :

\begin{lstlisting}[language=C]
result.allocated > SIZE_MAX / sizeof(element_type)
\end{lstlisting}

Parce que tester :

\[
allocated \times sizeof(T) > SIZE\_MAX
\]

pourrait déjà provoquer un overflow.
On transforme donc mathématiquement :

\[
x \times y > z
\]

en

\[
x > \frac{z}{y}
\]

pour éviter le débordement.

\subsection{5. Fonction vector\_push}

Convention choisie :
\begin{itemize}
    \item retourne le nouveau nombre d’éléments,
    \item retourne 0 en cas d’échec.
\end{itemize}

\begin{lstlisting}[language=C]
size_t vector_push(vector* v, element_type value) {
    if (v == NULL) {
        return 0;
    }

    // Tant que la capacite est insuffisante
    while (v->size >= v->allocated) {
        if (vector_enlarge(v) == NULL) {
            return 0;
        }
    }

    // insertion a la position size
    v->content[v->size] = value;

    v->size++;

    return v->size;
}
\end{lstlisting}

\subsection{6. Pourquoi un while et pas un if ?}

On aurait pu écrire :

\begin{lstlisting}[language=C]
if (v->size == v->allocated)
\end{lstlisting}

Mais utiliser :

\begin{lstlisting}[language=C]
while (v->size >= v->allocated)
\end{lstlisting}

rend le code plus général :

\begin{itemize}
    \item Permet d’ajouter plusieurs éléments d’un coup.
    \item Ne dépend pas d’une hypothèse « on ajoute toujours 1 élément ».
    \item Rend la fonction plus robuste.
\end{itemize}

Règle pratique :

\[
\boxed{\text{Si le corps modifie la condition, demandez-vous si un while n’est pas préférable à un if}}
\]

\subsection{7. Logique complète du push}

Deux cas :

\paragraph{Cas 1 : place disponible}

\begin{itemize}
    \item \texttt{size < allocated}
    \item on écrit directement à l’indice \texttt{size}
\end{itemize}

\paragraph{Cas 2 : plus de place}

\begin{itemize}
    \item appel à \texttt{vector\_enlarge}
    \item éventuellement plusieurs fois
    \item puis insertion
\end{itemize}

Les indices valides sont :

\[
0 \dots size - 1
\]

Le nouvel élément est inséré à :

\[
index = size
\]

Puis :

\[
size \leftarrow size + 1
\]

% chapter Pointers - 2 (end)
